\section{Mittelwert, Effektivwert und Leistung}

Diese Section behandelt die \textbf{zentralen Bewertungsgrössen} für periodische, nichtsinusförmige Spannungen und Ströme,
wie sie in Gleichrichtern, Stellern und Wechselrichtern auftreten.

\smallskip

\textbf{Zielbild:}
\begin{itemize}
    \item Definition von Mittelwert und Effektivwert
    \item Physikalische Interpretation (DC-Anteil, thermische Wirkung)
    \item Leistungsberechnung bei allgemeinen periodischen Signalen
    \item Anwendung auf stückweise definierte Zeitverläufe aus Übungen
\end{itemize}

\subsection{Periodische Signale und Grundbegriffe}

Ein Signal $x(t)$ heisst \textbf{periodisch} mit Periodendauer $T$, wenn gilt:
\[
x(t) = x(t+T) \qquad \forall t.
\]

Typische Beispiele aus der Leistungselektronik:
\begin{itemize}
    \item Gleichgerichtete Sinusspannung (Halb- und Vollwelle)
    \item Strom im R-L-Kreis bei Schaltbetrieb
    \item PWM-Ausgangsspannungen
\end{itemize}

\crd{\textrightarrow\ Mittelwert und Effektivwert werden immer über \textbf{eine volle Periode $T$} gebildet.}

\subsection{Zeitmittelwert (Gleichwert, DC-Anteil)}

\subsubsection{Definition}

Der \textbf{Zeitmittelwert} (Mittelwert) eines periodischen Signals $x(t)$ ist definiert als:
\[
\cbl{\overline{x} = \frac{1}{T} \int_{0}^{T} x(t)\,\diff t.}
\]

Interpretation:
\begin{itemize}
    \item Entspricht dem \textbf{Gleichanteil} des Signals.
    \item Bei Strom: mittlerer Transport von Ladung pro Zeit.
    \item Bei Spannung: DC-Anteil der Ausgangsspannung eines Gleichrichters.
\end{itemize}

\subsubsection{Typische Eigenschaften}

\begin{itemize}
    \item Symmetrische Wechselgrössen: $\overline{x}=0$.
    \item Gleichgerichtete Signale: $\overline{x}>0$.
    \item Für stückweise definierte Signale:
    \[
    \overline{x} = \frac{1}{T} \left( \int_{t_1}^{t_2} x_1(t)\,\diff t + \int_{t_2}^{t_3} x_2(t)\,\diff t + \dots \right).
    \]
\end{itemize}

\crd{\textrightarrow\ Mittelwert ist \textbf{nicht} gleich Effektivwert. Häufige Prüfungsfalle.}

\subsection{Effektivwert (RMS)}

\subsubsection{Definition}

Der \textbf{Effektivwert} eines periodischen Signals $x(t)$ ist definiert als:
\[
\cbl{x_{\mathrm{RMS}} = \sqrt{\frac{1}{T} \int_{0}^{T} x^2(t)\,\diff t}.}
\]

Interpretation:
\begin{itemize}
    \item Entspricht dem \textbf{Gleichwert}, der in einem Widerstand \textbf{die gleiche Verlustleistung} erzeugt.
    \item Thermische Wirkung.
    \item Grundlage für Dimensionierung von Leitern, Bauteilen und Verlusten.
\end{itemize}

\subsubsection{Vergleich Mittelwert vs. Effektivwert}

\begin{itemize}
    \item Mittelwert: lineare Mittelung.
    \item Effektivwert: quadratische Mittelung.
    \item Allgemein gilt:
    \[
    x_{\mathrm{RMS}} \ge |\overline{x}|.
    \]
\end{itemize}

\subsection{Momentanleistung und Wirkleistung}

\subsubsection{Momentanleistung}

Für allgemeine Spannungen und Ströme:
\[
\cbl{p(t) = u(t)\, i(t).}
\]

Diese Grösse schwankt zeitlich, auch im stationären Betrieb.

\subsubsection{Wirkleistung}

Die \textbf{mittlere Wirkleistung} ist definiert als:
\[
\cbl{P = \overline{p(t)} = \frac{1}{T} \int_{0}^{T} u(t)\,i(t)\,\diff t.}
\]

\subsubsection{Spezialfall: Ohmsche Last}

Für $u(t) = R\, i(t)$ gilt:
\[
\cbl{P = R \, I_{\mathrm{RMS}}^2 = \frac{U_{\mathrm{RMS}}^2}{R}.}
\]

Interpretation:
\begin{itemize}
    \item Nur der Effektivwert bestimmt die Verlustleistung.
    \item Unabhängig von der Signalform.
\end{itemize}

\subsection{Berechnung bei stückweise definierten Signalen}

In vielen Übungen ist $x(t)$ stückweise gegeben:
\[
x(t) =
\begin{cases}
x_1(t), & 0 < t < T_1,\\
x_2(t), & T_1 < t < T_2,\\
\dots
\end{cases}
\]

Dann gilt:

\paragraph{Mittelwert}
\[
\cbl{\overline{x} = \frac{1}{T} \sum_k \int_{T_k} x_k(t)\,\diff t.}
\]

\paragraph{Effektivwert}
\[
\cbl{x_{\mathrm{RMS}} = \sqrt{ \frac{1}{T} \sum_k \int_{T_k} x_k^2(t)\,\diff t }.}
\]

\paragraph{Wirkleistung}
\[
\cbl{P = \frac{1}{T} \sum_k \int_{T_k} u_k(t)\, i_k(t)\,\diff t.}
\]

\crd{\textrightarrow\ Bei stückweisen Verläufen immer korrekt die Zeitintervalle trennen.}

\subsection{Bezug zu Fourier-Reihen}

Für ein periodisches Signal mit Fourierkoeffizienten gilt (Parseval):

\[
\cbl{X_{\mathrm{RMS}}^2 = \frac{a_0^2}{2} + \sum_{n=1}^\infty \left( a_n^2 + b_n^2 \right).}
\]

Interpretation:
\begin{itemize}
    \item Effektivwert ist die Summe der Quadrate aller Harmonischen.
    \item Harmonische tragen direkt zur Verlustleistung bei.
\end{itemize}

\subsection{Skizze 1: Nichtsinusförmiger periodischer Strom}

\begin{center}
\begin{tikzpicture}[scale=1]
\draw[->] (0,0) -- (8,0) node[right] {$t$};
\draw[->] (0,0) -- (0,3) node[above] {$i(t)$};

\draw (0,1) -- (2,2.5) -- (4,1) -- (6,2.5) -- (8,1);

\draw[dashed] (0,1) -- (8,1) node[right] {$\overline{i}$};
\draw[dashed] (0,2) -- (8,2) node[right] {$i_{\mathrm{RMS}}$};

\draw (4,0) node[below] {$T/2$};
\draw (8,0) node[below] {$T$};
\end{tikzpicture}
\end{center}

\textbf{Aussage:}
\begin{itemize}
    \item Mittelwert ist die horizontale Mittel-Linie.
    \item Effektivwert liegt darüber.
\end{itemize}

\subsection{Skizze 2: Quadratbildung für Effektivwert}

\begin{center}
\begin{tikzpicture}[scale=1]
\draw[->] (0,0) -- (8,0) node[right] {$t$};
\draw[->] (0,0) -- (0,3) node[above] {$x^2(t)$};

\draw (0,0.5) -- (2,2.5) -- (4,0.5) -- (6,2.5) -- (8,0.5);

\draw[dashed] (0,1.5) -- (8,1.5) node[right] {$\overline{x^2}$};
\end{tikzpicture}
\end{center}

\textbf{Aussage:}
\begin{itemize}
    \item Zuerst quadrieren.
    \item Dann mitteln.
    \item Dann Wurzel ziehen.
\end{itemize}

\subsection{Skizze 3: Leistung bei ohmscher Last}

\begin{center}
\begin{tikzpicture}[scale=1]
\draw[->] (0,0) -- (8,0) node[right] {$t$};
\draw[->] (0,0) -- (0,3) node[above] {$p(t)$};

\draw (0,1.5) .. controls (1,2.5) .. (2,1.5)
      .. controls (3,0.5) .. (4,1.5)
      .. controls (5,2.5) .. (6,1.5)
      .. controls (7,0.5) .. (8,1.5);

\draw[dashed] (0,1.5) -- (8,1.5) node[right] {$P$};
\end{tikzpicture}
\end{center}

\textbf{Aussage:}
\begin{itemize}
    \item Momentanleistung schwankt.
    \item Wirkleistung ist der Mittelwert.
\end{itemize}

\subsection{Typische Fehlerquellen}

\begin{itemize}
    \item Mittelwert mit Effektivwert verwechselt.
    \item Quadrat beim RMS vergessen.
    \item Falsche Periodendauer $T$ eingesetzt.
    \item Stückweise Integrale nicht getrennt.
    \item Leistung mit $U_{\mathrm{RMS}}\cdot I_{\mathrm{RMS}}$ berechnet bei nichtohmscher Last.
\end{itemize}

\subsection{Minimal-Checkliste für Aufgaben}

\begin{enumerate}
    \item Periodendauer $T$ bestimmen.
    \item Signalform stückweise festlegen.
    \item Mittelwert berechnen.
    \item Effektivwert berechnen.
    \item Bei Leistung: $p(t)=u(t)i(t)$ bilden und mitteln.
    \item Plausibilität prüfen: $x_{\mathrm{RMS}} \ge |\overline{x}|$.
\end{enumerate}
